\chapter{Описание стенда и подготовка}
\label{ch:stand}

\section{Программная среда}

Работа выполняется в облачном редакторе \textbf{Figma}, доступном двумя
способами:
\begin{itemize}
  \item \textbf{Веб-версия} --- открыть \texttt{https://www.figma.com}
        в современном браузере (рекомендуется Chromium-семейство:
        Google~Chrome, Microsoft~Edge, Yandex~Browser);
  \item \textbf{Desktop-приложение} --- установить клиент для
        Windows/macOS со страницы загрузки. Функционал идентичен
        веб-версии; единственное преимущество --- работа в нативном окне
        и поддержка локальных шрифтов системы.
\end{itemize}
Для выполнения данной лабораторной работы достаточно браузерной версии.
Регистрация в Figma выполняется по электронной почте или через аккаунт
Google. Бесплатный тарифный план \texttt{Starter} позволяет создавать
до трёх файлов в одной команде, чего достаточно для учебного проекта.

% — Скриншот: общий вид рабочего пространства Figma после входа —
\begin{figure}[h]
  \centering
  \includegraphics[width=0.85\textwidth]{img/01_workspace}
  \caption{Рабочее пространство Figma после входа в аккаунт}
  \label{fig:workspace}
\end{figure}

\section{Подготовка рабочего пространства}

Поскольку макет включает пять смысловых страниц (главная, команда,
участник, турнирная таблица, магазин), а каждая страница --- ещё и три
адаптивные версии, удобно с самого начала разнести их по отдельным
страницам файла Figma.

Этапы:
\begin{enumerate}
  \item На главной странице Figma в команде \texttt{Drafts} нажать
        \texttt{New design file}. Откроется новый пустой файл.
  \item Переименовать файл в \texttt{Lab~11~--~PHYGITAL~FC}
        (двойной клик по названию вверху окна).
  \item На левой панели страница по умолчанию называется
        \texttt{Page~1}. Переименовать её в \texttt{Главная страница}.
  \item Кнопкой \texttt{+} в верхней части левой панели добавить ещё
        четыре страницы: \texttt{Команда}, \texttt{Участник},
        \texttt{Турнирная таблица}, \texttt{Магазин}.
\end{enumerate}

Между страницами файла удобно переключаться кликом по имени в левой
панели; все они шарят общую библиотеку компонентов, цветов и шрифтов
(см.\ раздел~\ref{ch:practice}).

\section{Создание десктоп-фрейма}

Десктоп-макет каждой страницы создаётся как фрейм фиксированной ширины.
На странице \texttt{Главная страница}:
\begin{enumerate}
  \item Выбрать инструмент \texttt{Frame}~(\texttt{F}).
  \item На правой панели в разделе \texttt{Design~\(\rightarrow\)~Desktop}
        выбрать \texttt{Desktop~--~1440\(\times\)1024} или вручную
        нарисовать прямоугольник нужного размера, затем растянуть его
        по высоте.
  \item Задать имя фрейма \texttt{Desktop~1440} (двойной клик по имени
        в левой панели слоёв).
  \item В разделе \texttt{Fill} задать линейный градиент-фон от
        \texttt{\#0B1024} (deep-navy) сверху к \texttt{\#11183A}
        (тёмно-синий) снизу --- основа выбранного стиля
        \emph{glassmorphism}.
\end{enumerate}

% — Скриншот: созданный фрейм Desktop 1440 на холсте —
\begin{figure}[h]
  \centering
  \includegraphics[width=0.85\textwidth]{img/02_frame_1440}
  \caption{Десктоп-фрейм \texttt{Desktop~1440} с градиентным фоном}
  \label{fig:frame_1440}
\end{figure}

\section{Импорт HTML-каркаса (плагин html.to.design)}

Чтобы не рисовать пятнадцать макетов вручную (\(5\) страниц
\(\times~3\) ширины), используется промежуточный HTML-каркас, который
автоматически импортируется во фреймы Figma:

\begin{enumerate}
  \item Открыть меню \texttt{Main~menu~\(\rightarrow\)~Plugins~\(\rightarrow\)
        Manage plugins}.
  \item В строке поиска ввести \texttt{html.to.design}; установить
        одноимённый плагин.
  \item Запустить плагин: \texttt{Plugins~\(\rightarrow\)
        html.to.design~\(\rightarrow\)~Run}.
  \item На вкладке \texttt{Upload file} загрузить
        \texttt{desktop/home.html}, в поле \texttt{Viewport~width}
        указать \texttt{1440}.
  \item Включить опции \texttt{Use~Autolayout},
        \texttt{Create~styles~\&~variables},
        \texttt{HTML~layer~names}; остальные опции выключить.
  \item Повторить для оставшихся четырнадцати HTML-файлов из
        \texttt{figma-mockup/}: каждый импорт уходит на свою страницу
        (\texttt{Команда}, \texttt{Участник}, \ldots) и попадает в
        фрейм нужной ширины.
\end{enumerate}

\section{Установка плагина Unsplash}

Плагин \texttt{Unsplash} даёт прямой доступ к каталогу бесплатных
фотографий из Figma. Он понадобится на этапе доводки макета, когда
плейсхолдеры из HTML-каркаса заменяются на тематические снимки
(стадион, болельщики, портреты игроков, мерч).

\begin{enumerate}
  \item Открыть меню \texttt{Main~menu~\(\rightarrow\)~Plugins~\(\rightarrow\)
        Manage plugins}.
  \item Перейти на вкладку \texttt{Browse plugins in Community}.
  \item В строке поиска ввести \texttt{Unsplash}; выбрать одноимённый
        плагин разработчика \textit{Unsplash}.
  \item Нажать \texttt{Install} (потребуется быть авторизованным
        в Figma).
\end{enumerate}

После установки плагин доступен через
\texttt{Resources~(Shift+I)~\(\rightarrow\)~Plugins~\(\rightarrow\)~Unsplash}
или контекстным меню по правому клику на выделенной фигуре:
\texttt{Plugins~\(\rightarrow\)~Unsplash}.

\section{Фирменный стиль чемпионата}

Перед началом проектирования зафиксирована визуальная система,
единая для всех пяти страниц:

\begin{itemize}
  \item \textbf{Палитра.} Фон --- линейный градиент от
        \texttt{\#0B1024} к \texttt{\#11183A}; основной акцент ---
        \texttt{\#3A8DFF} (electric-blue); дополнительный акцент ---
        \texttt{\#C7FF4F} (lime); поверхность карточек ---
        \texttt{rgba(255,255,255,0.08)} с обводкой
        \texttt{rgba(255,255,255,0.12)} и размытием
        \texttt{backdrop-filter:~blur(20px)} (стиль
        \emph{glassmorphism}).
  \item \textbf{Шрифты.}
        \texttt{Space~Grotesk} --- заголовки и крупные акцентные числа;
        \texttt{Inter} --- основной текст интерфейса;
        \texttt{JetBrains~Mono} --- табличная статистика и счёт матчей.
        Все три шрифта подключаются как Google~Fonts.
  \item \textbf{Иконография.} Phosphor~Icons либо Lucide --- линейные
        иконки толщиной \texttt{1.5\,px}, цвет \texttt{\#3A8DFF} либо
        белый \texttt{rgba(255,255,255,0.8)}.
  \item \textbf{Скругления.} Карточки --- \texttt{16\,px}; кнопки ---
        \texttt{12\,px}; чипы и теги --- \texttt{999\,px} (pill).
\end{itemize}

\endinput
